\documentclass[11pt,letterpaper]{article}
\usepackage[margin=1in]{geometry}
\usepackage{setspace}      % optional line

\usepackage[T1]{fontenc}
\usepackage[utf8]{inputenc}
\usepackage{lmodern}
\usepackage{microtype}
\usepackage{graphicx}
\usepackage{amsmath}
\usepackage{xcolor}
\usepackage{enumitem}

\begin{document}

\title{PINN Example: Solving the 1D Heat Equation}
\author{Cyrus Young}
\date{April 21, 2026}

\maketitle

\begin{figure}[ht]
    \centering
    \includegraphics[width=0.6\textwidth]{figures/PINN_img.png}
    \caption{PINN overview for the 1D heat equation.}
\end{figure}

\newpage

\subsection{How would one solve the 1D heat equation with a PINN?}

To make the PINN methodology concrete, consider the one-dimensional heat equation
\[
\frac{\partial T}{\partial t} = \alpha \frac{\partial^2 T}{\partial x^2} + q(x,t),
\]
defined on a spatial domain $x \in [0,L]$ and for $t > 0$, where $\alpha > 0$ is the thermal
diffusivity and $q(x,t)$ is a prescribed heat source. The problem is closed by an initial condition
\[
T(x,0) = f(x)
\]
and Dirichlet boundary conditions
\[
T(0,t) = g_0(t), \qquad T(L,t) = g_L(t).
\]
The goal is to find a function $T(x,t)$ satisfying all three of these requirements simultaneously.

Big picture: the implementation loop is $(x,t) \mapsto T_\theta(x,t)$ via a forward pass,
automatic differentiation to obtain $T_t$ and $T_{xx}$, evaluation of the PDE residual and
the total loss, gradients of that loss with respect to the weights and biases, and a parameter
update that produces the next set of predictions.

\subsubsection*{The neural network as a function approximator}

A PINN replaces the unknown solution $T(x,t)$ with the output of a feedforward neural network
$T_\theta(x,t)$. The network takes a single space-time point as its two-dimensional input and
produces a scalar temperature prediction:
\[
a^{[0]} = \begin{bmatrix} x \\ t \end{bmatrix}, \qquad a^{[N]} = T_\theta(x,t).
\]
The forward pass through the network proceeds layer by layer. For $\ell = 1, \dots, N$, a
pre-activation vector is formed by a linear transformation of the previous layer's output,
\[
z^{[\ell]} = W^{[\ell]}\, a^{[\ell-1]} + b^{[\ell]},
\]
and a nonlinear activation function $\sigma$ is then applied elementwise to produce the next
layer's output,
\[
a^{[\ell]} = \sigma\!\left(z^{[\ell]}\right), \qquad \ell = 1, \dots, N-1.
\]
Here $W^{[\ell]}$ is the weight matrix for layer $\ell$, $b^{[\ell]}$ is the corresponding bias
vector, $z^{[\ell]}$ is the pre-activation at layer $\ell$, and $a^{[\ell]}$ is its post-activation
output. The final layer $a^{[N]}$ is a single scalar and is interpreted directly as the predicted
temperature $T_\theta(x,t)$; no activation function is applied at the output.
The trainable parameters are $W^{[\ell]}$ and $b^{[\ell]}$ for $\ell=1,\dots,N$; the quantities
$z^{[\ell]}$, $a^{[\ell]}$, $T_\theta$, $T_t$, and $T_{xx}$ are intermediate values computed from the
current parameters during each forward pass.

\subsubsection*{Encoding the physics via automatic differentiation}

The network itself knows nothing about the heat equation. Physics enters through the loss
function. Because $T_\theta(x,t)$ is a differentiable composition of smooth functions, its
partial derivatives with respect to the inputs can be computed exactly by automatic
differentiation:
\[
T_t = \frac{\partial T_\theta}{\partial t}, \qquad
T_{xx} = \frac{\partial^2 T_\theta}{\partial x^2}.
\]
Substituting the network output into the heat equation yields the PDE residual
\[
R_\theta(x,t) = \frac{\partial T_\theta}{\partial t}
- \alpha\frac{\partial^2 T_\theta}{\partial x^2} - q(x,t).
\]
If $T_\theta$ were the true solution, this residual would vanish everywhere in the domain.
Training is designed to make it as small as possible.
Automatic differentiation is essential because the residual requires exact derivatives of the
network output with respect to its inputs at many points, which would be cumbersome and error-prone
to compute by hand or by finite differences.

\subsubsection*{The loss function}

The training objective is a weighted sum of three mean-squared loss terms, one for each
requirement imposed on the solution.

The \emph{PDE residual loss} penalises violations of the heat equation at a set of
$N_f$ interior collocation points, i.e., sampled space-time points where the PDE is enforced,
$\{(x_i^{(f)}, t_i^{(f)})\}$:
\[
\mathcal{L}_{\mathrm{PDE}}
= \frac{1}{N_f}\sum_{i=1}^{N_f}
\bigl|R_\theta(x_i^{(f)}, t_i^{(f)})\bigr|^2.
\]
The \emph{initial-condition loss} enforces the prescribed initial temperature profile
$f(x)$ at $N_0$ spatial points sampled at $t = 0$:
\[
\mathcal{L}_{\mathrm{IC}}
= \frac{1}{N_0}\sum_{i=1}^{N_0}
\bigl|T_\theta(x_i^{(0)}, 0) - f(x_i^{(0)})\bigr|^2.
\]
The \emph{boundary-condition loss} enforces the prescribed boundary temperatures $g_0$ and
$g_L$ at both ends of the domain, sampled at $N_b$ time points:
\[
\mathcal{L}_{\mathrm{BC}}
= \frac{1}{N_b}\sum_{i=1}^{N_b}
\bigl|T_\theta(0, t_i^{(b)}) - g_0(t_i^{(b)})\bigr|^2
+ \frac{1}{N_b}\sum_{i=1}^{N_b}
\bigl|T_\theta(L, t_i^{(b)}) - g_L(t_i^{(b)})\bigr|^2.
\]
These three terms are combined into the total loss
\[
\mathcal{L}
= \lambda_{\mathrm{PDE}}\,\mathcal{L}_{\mathrm{PDE}}
+ \lambda_{\mathrm{IC}}\,\mathcal{L}_{\mathrm{IC}}
+ \lambda_{\mathrm{BC}}\,\mathcal{L}_{\mathrm{BC}}.
\]

The integers $N_f$, $N_0$, and $N_b$ control how densely each constraint is sampled, and the scalar
coefficients $\lambda_{\mathrm{PDE}}$, $\lambda_{\mathrm{IC}}$, and $\lambda_{\mathrm{BC}}$ balance their relative
importance. The causal chain is: PDE, initial, and boundary errors contribute to $\mathcal{L}$;
$\mathcal{L}$ depends on $\theta$ through $T_\theta$ and its derivatives; gradients of $\mathcal{L}$
update $\theta$; and the updated $\theta$ changes the next forward pass, its derivatives, and the
residual.

At this point the total loss is a scalar function of all weights and biases,
\[
\theta = \{W^{[1]}, b^{[1]}, \dots, W^{[N]}, b^{[N]}\},
\]
so training computes the gradients
\[
\frac{\partial \mathcal{L}}{\partial W^{[\ell]}}, \qquad
\frac{\partial \mathcal{L}}{\partial b^{[\ell]}}
\]
for every layer and updates the parameters, for example by gradient descent,
\[
W^{[\ell]} \leftarrow W^{[\ell]} - \eta \frac{\partial \mathcal{L}}{\partial W^{[\ell]}},
\qquad
b^{[\ell]} \leftarrow b^{[\ell]} - \eta \frac{\partial \mathcal{L}}{\partial b^{[\ell]}}.
\]
The intermediate values $z^{[\ell]}$, $a^{[\ell]}$, $T_\theta$, $T_t$, and $T_{xx}$ are produced during
each forward pass and are not directly updated by the loss; instead, changing $W^{[\ell]}$ and
$b^{[\ell]}$ yields a new forward pass and therefore a new prediction $T_\theta(x,t)$. This is the
mechanism by which the PDE, initial, and boundary constraints influence the learned solution.

\subsubsection*{Summary of the workflow}

In brief, solving the 1D heat equation with a PINN proceeds as follows:
\begin{enumerate}[label=(\roman*)]
    \item Sample a space-time point $(x, t)$ in the domain (interior, initial line, or boundaries
          depending on the loss term being evaluated).
    \item Set the input activation
        \[
        a^{[0]} = \begin{bmatrix}x \\ t\end{bmatrix}.
        \]
    \item Propagate forward using
        \[
        z^{[\ell]} = W^{[\ell]} a^{[\ell-1]} + b^{[\ell]}, \qquad
        a^{[\ell]} = \sigma\!\left(z^{[\ell]}\right),
        \]
        and obtain $a^{[N]} = T_\theta(x,t)$.
    \item Compute the partial derivatives $\partial T_\theta / \partial t$ and
        $\partial^2 T_\theta / \partial x^2$ via automatic differentiation.
    \item Form the PDE residual $R_\theta(x,t)$ by substituting these derivatives into
        the heat equation.
    \item Aggregate $\mathcal{L}_{\mathrm{PDE}}$, $\mathcal{L}_{\mathrm{IC}}$, and $\mathcal{L}_{\mathrm{BC}}$ over
        their respective point sets to obtain $\mathcal{L}$.
    \item Back-propagate $\nabla_\theta \mathcal{L}$ and update $W^{[\ell]}$, $b^{[\ell]}$; repeat until
        $\mathcal{L}$ is acceptably small.
\end{enumerate}
At the end of training, $T_\theta(x,t)$ is a continuous, differentiable approximation of
the solution that satisfies the heat equation, initial condition, and boundary conditions
simultaneously, without an explicit mesh-based spatial discretisation or conventional time-marching scheme.

\subsubsection*{Implementation logic in one pass}

\begin{itemize}
    \item Fix the PDE, initial, and boundary data and choose sampling strategies for interior and boundary points.
    \item For each batch of points, run a forward pass to get $T_\theta$ and use automatic differentiation to get $T_t$ and $T_{xx}$.
    \item Compute $R_\theta$ at interior points and the IC/BC errors at their respective points.
    \item Form $\mathcal{L}$ as the weighted sum of these errors; $\mathcal{L}$ is a scalar function of $\theta$ only.
    \item Differentiate $\mathcal{L}$ with respect to all $W^{[\ell]}$ and $b^{[\ell]}$ and update them with the optimiser.
    \item Repeat until the loss and residuals reach acceptable levels.
\end{itemize}

\end{document}
